\subsection{多元自适应回归样条（MARS）}
\label{sec:regression-mars}

基函数回归需要人事先指定中心、节点与基类型；数据自己可能更知道该在哪里弯曲。\term{多元自适应回归样条}（multivariate adaptive regression splines, MARS）让基函数从数据中自动长出：以铰链函数为原料，按拟合改善逐步加入，再删除贡献小的项。Friedman在1991年提出这一数据驱动选基方法\scite{friedman1991mars}。

\subsubsection{模型结构}
\label{sec:regression-mars-structure}

一维铰链函数（hinge）为
\[
  h^+(x,t)=(x-t)_+=\max(0,x-t),
  \qquad
  h^-(x,t)=(t-x)_+=\max(0,t-x).
\]
MARS展开为
\begin{equation}
  f(\bm{x})=c_0+\sum_{m=1}^M c_m\prod_{k=1}^{K_m}h_{k}(\bm{x},t_{k}),
  \label{eq:regression-mars}
\end{equation}
乘积项表示变量之间的局部交互：单个铰链沿一个轴折点分段，多个铰链相乘才在矩形区域内共同起作用。

\subsubsection{参数学习}
\label{sec:regression-mars-learning}

\paragraph{前向选基}

从截距项开始，前向阶段加入成对铰链：每轮在所有候选铰链（每个变量、每个样本位置的两种方向）及其与已有项的乘积中，挑选使训练残差平方和下降最多的一项加入，并用最小二乘重估全部系数。结点、父基和样本支持的候选规则在此说明，候选项数量大，实际按网格或样本点限制。

\paragraph{后向删基与规模选择}

前向过程会把噪声也长出来，后向逐次删基形成规模序列，删除使广义交叉验证（GCV）最小的项。GCV把拟合残差按有效自由度打折，是留一误差的近似，免去为每个候选子集逐折重新拟合：
\begin{equation}
  \mathrm{GCV}=\frac{\mathrm{RSS}/n}{(1-C/n)^2},
  \label{eq:regression-mars-gcv}
\end{equation}
其中$C$为有效项数（含项数与交互阶数的调整）；在训练折内按验证或明确复杂度约定取GCV最小者。

\subsubsection{性质与使用条件}
\label{sec:regression-mars-theory}

由铰链及有限乘积的连续性说明整体连续，但结点未必可微；铰链在折点不可微，函数分段线性，变量乘积的交互阶数与单变量结点数量不是同一能力。固定基求系数与离散贪心选基的保证分开，GCV的搜索复杂度只是近似。外推到训练域外时项归零或线性延续，观测外铰链增长不保证合理，需要独立验证。与回归树对照：树的指示函数是硬划分，落入叶内用常数；MARS用连续铰链替代硬开关，在折点两侧平滑过渡。与固定基函数回归对照：两者都是基上的线性展开，差别在于基由人工指定还是由数据选择，后者把选基本身变成了学习问题。
